
\newcommand{\calC}{{\cal C}}

\section{Pascal's rule for computing $n$ choose $k$ (with long-form proof)}\label{dp: sec: n choose k}
\doHeader{Computing $n$ choose $k$}

Given non-negative integers $n$ and $k$,
the number of size-$k$ subsets of $\{1,2,\ldots, n\}$
is called \emph{$n$ choose $k$}, and denoted $n\choose k$.
Consider the problem of computing it:

\paragraph{{\prob{Computing~$n$ choose $k$}}}

\begin{mylist}
\item[\bf input:]
  non-negative integers $n$ and $k$

\item[\bf output:] $n\choose k$.
\end{mylist}

To design an algorithm, first we'll prove Pascal's rule:

\begin{lemma}[Pascal's rule]\label{dp: lemma: pascal}
  \(\displaystyle {n \choose k} =
  \begin{cases}
    1 & \text{ if } k = 0, \\
    0 & \text{ if } k > 0 \text{ and } n = 0,  \\
    {n-1 \choose k} + {n-1 \choose k-1} & \text{ otherwise.}
  \end{cases}
  \)
\end{lemma}
\begin{longFormProof}
  \step For $m,\ell\ge 0$, let $\calC(m,\ell)$ be the collection of all size-$\ell$ subsets of $\{1,2,\ldots,m\}$,
  so ${m \choose \ell} = |\calC(m, \ell)|$.

  \step If $k=0$, the empty set is the only size-$k$ subset of $\{1,2,\ldots,n\}$,
  so the rule holds in this case.

  \step If $k>0$ and $n=0$, there is no size-$k$ subset of $\{1,2,\ldots,n\}$,
  so the rule holds in this case. 

  \step So assume $n, k > 0$ and consider $\calC(n, k)$.  

  \step Each set in $\calC(n, k)$ either contains $n$ or it doesn't.

  \smallskip

  \step The sets in $\calC(n, k)$ that \emph{don't} contain $n$
  are just the sets in $\calC(n-1,k)$.

  \step By definition of $\calC$, there are $n-1 \choose k$ of these.

  \step\label{dp: step: choose 1}
  So, the number of sets in $\calC(n, k)$ that don't contain $n$
  is $n-1\choose k$.

  \smallskip

  \step To finish, we show that the number of sets in $\calC(n,k)$ that \emph{do} contain $n$ is $n-1\choose k-1$.

  \step For each set $S$ in $\calC(n,k)$ that does contain $n$,
  let $\phi(S) = S\setminus\{n\}$ be the set with $n$ removed.

  \step Then $\phi$ is a bijection (a one-to-one mapping)
  between the sets in $\calC(n, k)$ that contain $n$
  and the sets in $\calC(n-1, k-1)$.
  (That is, $\phi(S)$ is in $\calC(n-1, k-1)$,
  and, for each $S'\in\calC(n-1,k-1)$,
  there is exactly one set $S\in\calC(n,k)$ that contains $n$
  such that $\phi(S)=S'$---namely $S = S'\cup\{n\}$.)

  \step Because there is such a bijection,
  the number of sets in $\calC(n, k)$ that contain $n$
  must equal the number of sets in $\calC(n-1, k-1)$.
  By definition this number is $n-1\choose k-1$.

  % [review 2026-09-27] fix: the last step concluded C(n-1,k) + C(n,k); the second term (sets containing n) is C(n-1,k-1)
  \step By this and Step~\ref{dp: step: choose 1}, the size of $\calC(n,k)$
  is ${n-1\choose k} + {n-1 \choose k-1}$.
  
\end{longFormProof}

The proof of Pascal's rule illustrates the first main principle of dynamic programming.
Namely, to reason about a collection of objects
(in this case the size-$k$ subsets of $\{1,\ldots,n\}$)
the proof partitions the collection into a few types---in this case two:
\begin{mylist}
\item[(A)] Those subsets not containing the element $n$.
\item[(B)] Those subsets containing the element $n$.
\end{mylist}
Then it handles each type
by noting that, restricted to objects of that particular type,
the problem at hand is somehow equivalent to a smaller instance of the original problem.

In the current example, the proof observes that the sets of type (A) are just the size-$k$ subsets of $\{1,\ldots,n-1\}$,
so counting the type-(A) sets is equivalent to the problem of counting the size-$k$ subsets of $\{1,\ldots,n-1\}$.
This is a smaller instance of the original problem.

Similarly, the sets of type (B) correspond bijectively (via the deletion of element $n$)
to the size-$(k-1)$ subsets of $\{1,\ldots,n-1\}$.
So counting the type-(B) sets is equivalent to counting the size-$(k-1)$ subsets of $\{1,\ldots,n-1\}$.
Again, this is a smaller instance of the original problem.

This understanding of the recursive structure of the objects in question (subsets)
yields a recurrence relation (in this case, Pascal's rule).
This in turn yields a recursive, divide-and-conquer algorithm for the problem.
In this case, for computing $n \choose k$:

\smallskip

\begin{myframed}
  \begin{steps}
  \item[{\textsf{countSubsets}$(n, k)$:}] \comment{--- $n, k \ge 0$}
    \step return
    $\begin{cases} 1 & \text{ if } k=0, \\
      0 & \text{ if } k > 0 \text{ and } n = 0,\\
      \textsf{countSubsets}(n-1, k) + \textsf{countSubsets}(n-1, k-1) & \text{ otherwise.}
    \end{cases}$
  \end{steps}
\end{myframed}

Correctness of the algorithm follows by induction and Lemma~\ref{dp: lemma: pascal}.

This seems like a classic \emph{divide-and-conquer} approach.
But beware!  This implementation is not efficient---indeed, its run time is exponential!  Consider its recursion tree,
which has branching factor two, and depth between $k$ and $n$,
so somewhere between $2^k$ and $2^n$ nodes!

There are two standard work-arounds for this obstacle.
We'll discuss them in the next example.

\subsection*{Exercises}

\exercise\emph{(example with solution)}\label{dp: exercise: choose}
The proof of Lemma~\ref{dp: lemma: pascal} uses a bijection $\phi$ between
those sets in $\calC(n, k)$ that contain $n$ and the sets in $\calC(n-1, k-1)$.
Show this bijection for $n=5$ and $k=3$.
For reference, here is all of $\calC(5,3)$, and its subset $\calC(4,3)$:
\[
  \calC(5,3)\left\{
  \begin{array}{l}
    \left.\begin{array}{@{\!}l}
            \{1,2,3\} \\
            % [review 2026-09-27] fix: the set {1,2,4} was missing, so the display showed only 9 of the 10 sets of C(5,3) and 3 of the 4 sets of C(4,3); added it
            \{1,2,4\} \\
            \{1,3,4\} \\
            \{2,3,4\}
          \end{array}
    \right\} \calC(4,3) \\
    \{1,2,5\} \\
    \{1,3,5\} \\
    \{1,4,5\} \\
    \{2,3,5\} \\
    \{2,4,5\} \\
    \{3,4,5\}
  \end{array}
  \right.
\]

\solution
\[
  \begin{array}{ccc}
    \footnotesize\shortstack{sets in $\calC(5,3)$ \\ that contain $5$}
    & \shortstack{$\phi$\\$\longleftrightarrow$}
    & \footnotesize\shortstack{sets in $\calC(4,2)$} \\ \hline
    \{1,2,5\} && \{1, 2\} \\
    \{1,3,5\} && \{1, 3\} \\
    \{1,4,5\} && \{1, 4\} \\
    \{2,3,5\} && \{2, 3\}\\
    \{2,4,5\} && \{2, 4\}\\
    \{3,4,5\} && \{3, 4\}
  \end{array}
\]
\subsection*{External sources}
\begin{itemize}
\item \url{https://en.wikipedia.org/wiki/Pascal\%27s_rule}
\end{itemize}

%%% Local Variables:
%%% mode: latex
%%% TeX-master: "dynamic_programming"
%%% End:
